Teniendo en cuenta la descripción del sistema, se plantean los requisitos a satisfacer por parte de la aplicación. Estos constituyen la base sobre la cual se desarrolla el sistema y permiten definir de manera clara qué funcionalidades se esperan y qué condiciones deben cumplir. Bajo esta premisa se diferencian los requisitos funcionales, que detallan las operativas disponibles para los usuarios, como la interacción con el tutor socrático, la moderación con IA y la ejecución de código. Por su parte, los requisitos no funcionales establecen las restricciones técnicas, de seguridad y de rendimiento exigidas a la arquitectura.